\section{Overview}%
\label{sec:chap1Overview}
\autoBookmark{sec}{Overview}
Distributed Acoustic Sensing (DAS) converts a standard optical fibre into an array of thousands of virtual vibration sensors. An interrogator that analyses the light backscattered from the glass can monitor tens of kilometres of cable in real time, which makes the technology attractive for pipeline surveillance, railway monitoring, structural health monitoring, and perimeter security. Detecting that a disturbance has occurred is of limited operational value on its own. An installation becomes useful when it identifies the event that caused the disturbance, whether a vehicle, a person at a fence, or an excavator near a buried cable, and when that identification remains reliable under the conditions of an actual deployment.

This thesis develops and evaluates a lightweight multi-domain machine learning pipeline for event classification in DAS data. Each analysis window is described by named descriptors from four signal domains, namely time, frequency, time--frequency, and space; the descriptors are fused into a single vector, reduced by consensus feature selection, and passed to classifiers that train and predict on the central processing unit (CPU) of an embedded board. The pipeline was designed on the nine-class perimeter-security dataset of Tomasov \etal~\cite{tomasov2025dataset}, the primary dataset of the thesis, and then applied without modification to the six-class dataset of Cao \etal~\cite{caoxin2023opendataset}, which was recorded with an unrelated sensing system. Extending the evaluation to the second dataset required a precise definition of an independent sample in DAS data, and that definition showed that the evaluation protocols used with both datasets allow information to pass from the training data to the test data. Every result in this thesis is therefore obtained on partitions whose separation of training and test data has been verified in code.

The chapter is organised as follows. \refs{sec:ResearchBackground} describes the operating principle of DAS and the difficulties of event classification, and \refs{sec:ResearchProblemStatement} states the problem that the thesis addresses. \refs{sec:ResearchQuestions} and \refs{sec:ResearchObjectives} present the research questions and objectives, \refs{sec:SignificanceOfTheStudy} and \refs{sec:ResearchScope} the significance and scope of the study, and \refs{sec:ThesisContribution} its contributions and their novelty. \refs{sec:ThesisStructure} outlines the remaining chapters.

\section{Research Background}
\label{sec:ResearchBackground}
\autoBookmark{sec}{Research Background}
A DAS interrogator launches coherent laser pulses into the fibre and analyses the Rayleigh backscatter produced by refractive-index inhomogeneities along the glass, a technique known as phase-sensitive optical time-domain reflectometry ($\Phi$-OTDR)~\cite{he2021review,liu2025advances}. Acoustic and seismic disturbances near the fibre strain the glass and modulate the phase of the backscattered light, so that a vibration can be both localised and quantified at any point along the cable~\cite{hartog2017introduction}. Because a single interrogator covers tens of kilometres at metre-scale spatial resolution, DAS has become a major research frontier in acoustic sensing~\cite{shang2022das}.

Deployments now span pipeline surveillance, railway protection, road-traffic monitoring, and campus or border intrusion detection~\cite{wu2024pipeline,yurekli2025railway,deng2025roadtraffic,juarez2005intrusion}. In perimeter security the task is one of classification rather than detection: an operator must establish whether a disturbance originates from a passing car, a person climbing a fence, an excavator working near the cable, or routine background activity~\cite{tomasov2025perimeter}. The assigned label determines whether an alarm is raised, and the practical value of an installation therefore depends largely on classification accuracy.

DAS measurements are commonly represented as distance--time maps or time--frequency matrices from which descriptors are derived~\cite{kandamali2022ml}. Early systems relied on handcrafted signal processing: band-pass filtering was followed by the extraction of energy, spectral, and time--frequency features, which were supplied to machine learning models such as support vector machines and random forests~\cite{bublin2021event,breiman2001randomforests}. More recent work has shifted to supervised deep learning. Convolutional neural networks (CNNs) predominate, and recurrent, transformer, graph, and semi-supervised architectures have been proposed to capture long-range spatio-temporal structure or to reduce the labelling effort~\cite{tomasov2023enhancing,li2020convlstm,dione2025intrusion,jeong2026s2dygnn,kang2024semisupervised}.

Accuracy alone does not make a system deployable. Delayed or incorrect predictions slow the response to genuine events, false alarms consume operator time and erode trust, and missed detections allow intrusion, tampering, or nearby construction to pass unnoticed~\cite{han2024intrusion}. A practical classifier must therefore also be robust, fast enough for the data rate of the interrogator, interpretable to the personnel who act on its alerts, and reliable on every class rather than only on average~\cite{deng2025roadtraffic,yurekli2025railway}.

Such a classifier typically runs on an embedded computer beside the interrogator, so it must also be lightweight. The term is often used without a precise definition. In this thesis it refers to four measured quantities: the stored size of the trained model, the time needed to classify one window, the total processing time per window including feature extraction, and the time and hardware needed for training. A pipeline is regarded as lightweight when its model is small enough to be stored on and transferred to an edge device, when it trains within minutes on CPU cores without a graphics processing unit (GPU), when it classifies a window in well under a millisecond, and when its total processing time per window fits within the interval at which new windows arrive, using only a few CPU cores of an embedded board. The four quantities do not necessarily agree. A small network may still require hours of GPU training, and a classifier that predicts within microseconds may depend on a feature extractor that dominates the processing time. In a pipeline built on engineered features, the number of descriptors enters every one of these quantities, because each descriptor must be computed for every window, held in memory, and searched during training.

\section{Research Problem Statement}
\label{sec:ResearchProblemStatement}
\autoBookmark{sec}{Research Problem Statement}
The first difficulty lies in the signal and its representation. Each DAS channel produces a long, high-rate time series contaminated by ambient noise, coherent fading, and environmental interference, which complicates feature extraction and raises the cost of real-time processing~\cite{huang2024noise}. The signature of a given event also varies considerably with its distance from the fibre and with the quality of the coupling between the fibre and the ground. A useful representation must capture the characteristics of the event rather than those of the conditions under which it was recorded. Earlier studies have described DAS events with statistics of the waveform, spectral and cepstral coefficients, spectrogram and wavelet representations, and patterns across neighbouring channels, but most combine only two or three of these domains (Chapter~\ref{chap:LR}), and the contribution of each domain has seldom been measured on the same data.

The second difficulty is a trade-off between two modelling approaches. Engineered feature sets are inexpensive to evaluate and interpretable, but they are often large and redundant, and in the small, imbalanced datasets typical of DAS they are prone to overfitting unless their dimensionality is reduced in a principled manner~\cite{wani2025dimensionality,guyon2003variable}. Deep networks attain the highest reported accuracies, but they generally require GPUs, contain many parameters, offer limited insight into their decisions, and remain vulnerable to dataset shift~\cite{shahabudin2025review,shi2025benchmarking}. Neither profile is well suited to the resource-constrained, audit-sensitive setting of an edge-deployed monitoring system~\cite{deng2025roadtraffic}.

The third difficulty concerns evaluation: the protocols used in this field may permit information to pass between the training and test data. Two properties of DAS data make naive partitioning unsafe. First, event datasets are constructed by sliding an analysis window along a continuous stream. When the hop is shorter than the window, neighbouring samples share raw measurements, and a randomly selected test window may be three-quarters identical to a training window. Second, a recording session fixes the fibre state, the coupling, the interrogator settings, the weather, and often the operator, so that the samples of one session are repeated observations of a single physical situation rather than independent draws. A sufficiently flexible model can then achieve a high score by recognising the session rather than the event. Both properties are present in the two datasets examined in this thesis (Chapter~\ref{chap:ResultsAndDiscussion}). An accuracy measured under such conditions is an upper bound whose relation to deployment performance is unknown, and it confounds the merit of a method with the leniency of its protocol.

These difficulties have mostly been studied separately. Feature engineering, multi-domain representations, and deep architectures have each received attention, but few studies combine descriptors from all four signal domains, reduce them by an explicit selection procedure, and handle class imbalance within a pipeline light enough for an embedded board, and fewer still compare such a pipeline with recent deep learning models on identical partitions of more than one dataset. Where computational cost is reported, it is usually the inference time of the classifier, and rarely the stored size of the model, its training cost, or the cost of computing its input. The problem addressed by this thesis therefore has three parts: which feature representations earlier studies have used for DAS event classification and which domains they leave uncovered; how descriptors from all four domains can be fused, reduced, and balanced in a lightweight pipeline; and how such a pipeline performs against published deep learning models when both are evaluated on the same partitions, free of leakage.

\section{Research Questions}
\label{sec:ResearchQuestions}
\autoBookmark{sec}{Research Questions}
The research questions follow the design cycle. RQ1 examines the feature representations used in earlier work, RQ2 the development of the pipeline, and RQ3 its evaluation. Each question addresses one of the research gaps identified in \refs{sec:chap2ResearchGap}:
\begin{enumerate}[label=\textbf{RQ\arabic*.}, leftmargin=*, itemsep=4pt]
	\item What feature representations are used to classify events in DAS data?
	\item How can multi-domain feature fusion, consensus feature selection and class balancing be combined in a lightweight machine learning pipeline for DAS event classification?
	\item How does the lightweight multi-domain machine learning pipeline perform compared with published deep learning models for DAS event classification?
\end{enumerate}

\section{Research Objectives}
\label{sec:ResearchObjectives}
\autoBookmark{sec}{Research Objectives}
Each research question has a corresponding research objective (RO), and each objective has one main action:
\begin{enumerate}[label=\textbf{RO\arabic*.}, leftmargin=*, itemsep=4pt]
	\item To analyse the feature representations used for event classification in DAS data.
	\item To develop a lightweight multi-domain machine learning pipeline for event classification in DAS data.
	\item To evaluate the lightweight multi-domain machine learning pipeline for event classification in DAS data.
\end{enumerate}
RO1 examines the descriptors used in earlier work, identifies the domains that they leave uncovered, and measures on both datasets how well each domain classifies events on its own. RO2 designs a pipeline that covers all four domains while remaining lightweight. RO3 evaluates that pipeline against published deep learning models on identical partitions, in terms of classification performance, generalisation to new recordings, model size, and processing time. Chapter~\ref{chap:methodology} maps each pair of question and objective to an activity of the Design Science Research Methodology (DSRM) of Peffers et al.~\cite{peffers2007dsrm} (Table~\ref{tab:dsrm_mapping}), and Chapter~\ref{chap:Conclusion} assesses whether each objective has been achieved.

\section{Significance of the Study}
\label{sec:SignificanceOfTheStudy}
\autoBookmark{sec}{Significance of the Study}
The study is relevant to two groups: the people who must act on DAS alerts at a monitored site, and the researchers who build and compare DAS classifiers.

\subsection{Practical Significance}
\label{sec:PracticalApplications}
\autoBookmark{subsec}{Practical Significance}
Perimeter events must be recognised by those responsible for the sites that a fibre protects. They include security operators at campuses and other fenced facilities~\cite{tomasov2025perimeter}, the operators of pipelines, railways, and road networks monitored with DAS~\cite{wu2024pipeline,yurekli2025railway,deng2025roadtraffic}, and the owners of buried cables, for whom excavation near the route is a direct threat~\cite{bublin2021event}. For these users, a detected vibration is useful only when it carries a label that tells them whether to respond.

The volume of data explains why the classification is best performed on an embedded device beside the interrogator. The interrogator of the primary dataset samples $1{,}700$ positions along the fibre at 20\,kHz, which amounts to 34 million samples per second. Classifying the windows at the site means that only the label of an event, its position and time, and the descriptors that supported it need to leave the site, rather than the raw stream. Such a design reduces the bandwidth required for remote monitoring, permits operation where the network link is limited, and avoids the delay of sending data to a central server before an alert can be raised, which are the constraints that operational DAS systems must meet~\cite{deng2025roadtraffic}. A classifier whose learning stage runs on CPU cores can also be retrained at the site when conditions change, for example after a fence is rebuilt or the mix of nearby activity shifts, without a GPU. Because each decision rests on named descriptors, an operator or an auditor can see which properties of the signal produced an alert, which supports appropriate trust in automated alerting.

\subsection{Research Significance}
\label{sec:TheoreticalAdvancements}
\autoBookmark{subsec}{Research Significance}
For researchers, the study provides three resources. The first is a review of the feature representations used in twelve DAS event-classification studies that records, for each study, the signal domains covered and the protocol under which its accuracy was obtained; it shows which descriptors have been tried and which combinations have not. The second is a pipeline design in which every descriptor is named, assigned to one of four domains, defined by an equation, and traceable to the code that computes it, so that the design can be reproduced, modified, and inspected. The third is an evaluation on two unrelated public datasets in which the pipeline and published deep learning models are trained on identical partitions and measured on one embedded board, together with the partition manifests and verification code that allow other researchers to evaluate their own models under the same conditions.

\section{Research Scope}
\label{sec:ResearchScope}
\autoBookmark{sec}{Research Scope}
The study uses two public DAS datasets with clearly separated roles. The primary dataset is the nine-class perimeter-security corpus of Tomasov \etal~\cite{tomasov2025dataset}, recorded on a university campus with a phase-sensitive OptaSense ODH-F interrogator; the pipeline was designed on this dataset, and the principal results are reported on it. The evaluation dataset is the six-class intensity-based $\Phi$-OTDR release of Cao \etal~\cite{caoxin2023opendataset}. It is subjected to the identical process, from sample manifests and leakage audit to partitioning, preprocessing, feature extraction, selection, balancing, classification, and diagnostics, with the same settings and without per-dataset tuning. The two acquisitions share no hardware, sensing principle, channel geometry, or event taxonomy, so a conclusion that holds on both can be attributed to the method rather than to a single site. Four event types nevertheless occur in both taxonomies, which permits a direct test of whether a model trained on one system transfers to the other.

Methodologically, the study covers handcrafted multi-domain feature extraction, ensemble feature selection, class balancing, lightweight machine learning classifiers (linear, single-tree, and tree-ensemble models), and the design and verification of a leakage-controlled evaluation protocol. The CNN published with each dataset and five architectures published from 2024 onwards (S2-DyGNN, DASViT1D, DAStatFormer, and graph attention and graph convolutional baselines) are reimplemented from their source descriptions and trained on identical partitions to serve as comparators. The development or exhaustive tuning of new deep architectures lies outside the scope of the study. Embedded execution is the target deployment setting: all runtimes are measured on an NVIDIA Jetson AGX Orin, and every model, comprising the pipeline and each baseline, is described on that platform by its stored size, its classification time per window, its total processing time per window, and its training time.

Robustness to weather and seasonal variation, deployment on further interrogator systems, scaling to the full channel count of an interrogator, and integration with alarm-management software are reserved for future work. Generalisation across recordings, acquisition campaigns, and sensing systems is examined within the limits of the available data, and negative outcomes are reported as measured.

\section{Thesis Contribution}
\label{sec:ThesisContribution}
\autoBookmark{sec}{Thesis Contribution}
The thesis makes three contributions, one for each research objective. Chapter~\ref{chap:methodology} describes the methods underlying each contribution, and Chapter~\ref{chap:ResultsAndDiscussion} reports the corresponding results:
\begin{enumerate}[label=\textbf{C\arabic*.}, leftmargin=*, itemsep=4pt]
	\item \textbf{A review of the feature representations used for event classification in DAS data} (RO1; Chapter~\ref{chap:LR}). For twelve studies, the review records the descriptors used, the signal domains they cover, the model, the reported accuracy, and the evaluation protocol under which that accuracy was obtained. The review establishes that earlier studies drew their descriptors from one to three of the four domains, most often from two; that none combined time, frequency, time--frequency, and spatial descriptors; and that none reduced its feature set by agreement among several selection criteria.
	\item \textbf{A lightweight multi-domain machine learning pipeline for event classification in DAS data} (RO2; Chapter~\ref{chap:methodology}). The pipeline describes each window with $1{,}002$ named descriptors from the time (112), frequency (90), time--frequency (739), and spatial (61) domains and fuses them into one vector in which every descriptor keeps a domain prefix. A unanimous vote of eight feature selectors from the filter, wrapper, embedded, and attribution families reduces that vector, class imbalance is handled within the training data by a balancing strategy chosen on validation data, and seven lightweight classifiers train and predict on CPU cores without a GPU. The extractor, the selection ensemble, and the training procedure were implemented for this study as a Python library. The design was developed on the Tomasov dataset and applied unchanged to the Cao dataset.
	\item \textbf{An evaluation of the pipeline against published deep learning models on two unrelated public datasets} (RO3; Chapter~\ref{chap:ResultsAndDiscussion}). The CNN published with each dataset and five graph and transformer architectures published from 2024 onwards are trained on the same partitions as the pipeline, partitions whose separation of training and test data is verified in code before any model is trained (\refss{sec:chap3Phase1}). The evaluation reports per-class classification performance; generalisation to unseen recordings, unseen acquisition dates, and a second sensing system; and the model size, training time, and processing time per window of every model on one embedded board.
\end{enumerate}

\subsection{Novelty of the Work}
\label{sec:ThesisNovelty}
\autoBookmark{subsec}{Novelty of the Work}
The novelty of this work lies in its feature engineering rather than in a new learning algorithm. The transforms, selectors, resamplers, and classifiers used by the pipeline are established; the contribution is the way in which they are combined and implemented for DAS data. The feature engineering carried out in this study comprises five elements: the selection and implementation of $1{,}002$ descriptors across four domains, including band ratios and indices matched to the frequency bands of DAS events and spatio-temporal descriptors that follow the energy of an event across neighbouring channels; the choice of a short-time Fourier transform resolution that gives the spectrogram texture descriptors enough frames to be computed; the replacement of descriptors that per-window normalisation leaves undefined, such as the coefficient of variation, by scale-invariant equivalents; the fusion of the four domain matrices with a domain prefix on every descriptor name; and the reduction of the fused set by unanimous consensus of eight selectors fitted on training data only. Chapter~\ref{chap:methodology} defines each element and shows the code that implements the principal steps.

Interpretability follows from the same design. Because every descriptor has a name, a domain, and a formula, a reader can see which descriptors a fitted model relies on and check each one against its definition. For example, \texttt{TF:flux\_max} is the largest change in spectrogram magnitude between consecutive frames of a window, so an alert that relies on it can be traced to an abrupt change in the spectrum. \refss{sec:resWhiteBox} shows that this kind of inspection exposed descriptors that carried information about the recording rather than the event, a problem that the learned representation of a network would not have revealed. Table~\ref{tab:novelty} compares the four elements of the design with the prevailing practice in the DAS studies reviewed in Chapter~\ref{chap:LR}, and Chapter~\ref{chap:Conclusion} revisits them with the measured evidence for each.

\begin{table}[!htbp]
	\centering
	\caption{Elements of the feature engineering introduced in this thesis, compared with the prevailing practice in the reviewed DAS studies.}
	\label{tab:novelty}
	\footnotesize
	\setlength{\tabcolsep}{4pt}
	\renewcommand{\arraystretch}{1.2}
	\begin{tabularx}{\textwidth}{@{}>{\raggedright\arraybackslash}p{3.0cm} >{\raggedright\arraybackslash}X >{\raggedright\arraybackslash}X@{}}
		\toprule
		\textbf{Element} & \textbf{Prevailing practice} & \textbf{Introduced in this thesis}\\
		\midrule
		Coverage of the signal domains & Descriptors from one to three domains; none of the twelve reviewed studies covers all four (Table~\ref{tab:feature-domain-coverage}) & $1{,}002$ descriptors from the time, frequency, time--frequency, and spatial domains, computed from one preprocessed window in a single pass\\
		\addlinespace[3pt]
		Descriptors designed for DAS windows & Spatial information enters as a raw time--space matrix or channel graph for a network, or as statistics computed per channel & Spatio-temporal descriptors that track the energy centroid of an event across channels within the window; texture descriptors of the quantised log-spectrogram; band ratios and tonal, transient, and cyclostationary indices\\
		\addlinespace[3pt]
		Reduction of the feature set & No reduction, or a single method such as principal component analysis & Unanimous vote of eight selectors from the filter, wrapper, embedded, and attribution families, all fitted on training rows only\\
		\addlinespace[3pt]
		Traceability of the inputs & Learned representations or unnamed coefficient vectors, explained after training if at all & Every descriptor carries a name, a domain prefix, an equation, and a location in the code, so a decision can be checked against the definitions of the descriptors it used\\
		\bottomrule
	\end{tabularx}
\end{table}

\section{Thesis Structure}
\label{sec:ThesisStructure}
\autoBookmark{sec}{Thesis Structure}
The remaining four chapters follow the activities of the DSRM, which Chapter~\ref{chap:methodology} describes and \reff{fig:dsrm_flow} summarises: the review of prior work, the design of the pipeline, its demonstration and evaluation, and the consolidation of the findings. Each chapter opens with an overview of its structure and closes with a summary that relates its content to the research questions.
\begin{itemize}[leftmargin=*, itemsep=4pt]
	\item \textbf{Chapter~\ref{chap:LR} (Literature Review)} reviews the relationship between artificial intelligence, machine learning, and deep learning; the stages of a signal-classification pipeline, from feature extraction, combination, and selection to class-imbalance handling and modelling; the principles of DAS and its signal processing; machine learning, deep learning, and graph approaches to DAS event classification; deployment requirements; and evaluation protocols, including the leakage controls on which this study relies. It tabulates the feature representations of twelve prior studies together with the evaluation protocol of each, answers the literature part of RQ1, and identifies the three research gaps that motivate RQ1 to RQ3.

	\item \textbf{Chapter~\ref{chap:methodology} (Methodology)} justifies the adoption of the DSRM, maps the research questions and objectives to its activities, and describes the pipeline in detail: the primary and evaluation datasets and their sample manifests, the leakage found in their released protocols and the group-aware partitions that replace them, the preprocessing chain, four-domain feature extraction and fusion, consensus feature selection, class balancing, classifier configuration, the baseline models, and the procedures used to evaluate them.

	\item \textbf{Chapter~\ref{chap:ResultsAndDiscussion} (Results and Discussion)} reports the results in the order of the methods, for the Tomasov dataset first and the Cao dataset second: the leakage audit and the realised partitions; classification performance by feature domain, which completes the answer to RQ1; feature selection, class balancing, and the comparison of classifiers, which answer RQ2; and the per-class results, statistical tests, spatial-domain analysis, generalisation tests, cross-system transfer, and the measured model size, training time, and processing time per window. The discussion then compares the two datasets, compares the pipeline with the published CNNs and the recent architectures on identical partitions, which answers RQ3, and contrasts the white-box character of the pipeline with the black-box character of the deep baselines.

	\item \textbf{Chapter~\ref{chap:Conclusion} (Conclusions and Future Directions)} revisits each research objective, states the contributions of the work with the evidence for each, presents the limitations, and relates each limitation to a direction for future work.
\end{itemize}
